%% GENERATED by tools/gen_latex.py -- edit content/ch16.py instead.
%% Build:  latexmk -pdf ch16.tex     (or: pdflatex ch16.tex, twice)
\documentclass[aspectratio=169,11pt,t]{beamer}
\usepackage{itbpro}

\coursetitle{Text Generation}
\coursesubtitle{Building a GPT from the parts we already have, then fine-tuning a billion-parameter model into a chatbot --- and being honest about what none of it fixes.}
\coursekicker{Chapter 16}
\coursesource{Chollet \& Watson, Deep Learning with Python 3e --- chapter 16}
\courseduration{4 hours (3 sessions)}
\coursepresenter{Prof. Bambang Riyanto Trilaksono $\cdot$ Rahman Indra Kesuma, S.Kom., M.Cs.}{Lead Instructor $\cdot$ Teaching Assistant}
\coursebrand{AI for Professional \textperiodcentered\ ITB \textperiodcentered\ Ch 16}

\title{Text Generation}
\author{Prof. Bambang Riyanto Trilaksono $\cdot$ Rahman Indra Kesuma, S.Kom., M.Cs.}
\date{}

\begin{document}
\covertitle

\framekicker{Learning outcomes}
\begin{frame}[shrink=25]{Session Objectives}
\leadpar{By the end of this session, participants can:}
\begin{isteps}
\item Trace the line from \textbf{LSTM in 1997 to ChatGPT}, and name the three ingredients GPT combined.
\item Explain why GPT is \textbf{decoder-only}, and what that costs and buys.
\item Build a \textbf{mini-GPT}: data pipeline, tied embeddings, warmup schedule, and a logits-based loss.
\item Implement and compare \textbf{greedy, random, temperature, and top-K} sampling, and say what each failure mode looks like.
\item Explain \textbf{key-value caching} and why generation is slow without it.
\item \textbf{Instruction fine-tune} a pretrained LLM, and use \textbf{LoRA} to make it fit in accelerator memory.
\item Describe \textbf{RLHF, multimodal input, RAG, and chain-of-thought training} well enough to judge which one a given problem needs.
\item State the \textbf{structural limits}: hallucination, prompt sensitivity, energy cost, and a pretraining-data supply that is running out.
\end{isteps}

\vskip 6pt
\vskip 4pt
\reslink{Book}{Chapter 16 --- full text}{https://deeplearningwithpython.io/chapters/chapter16\_text-generation}
\reslink{Notebook}{01\_mini\_gpt\_data\_pipeline.ipynb}{../../notebooks/ch16/01\_mini\_gpt\_data\_pipeline.ipynb}
\reslink{Notebook}{02\_training\_mini\_gpt.ipynb}{../../notebooks/ch16/02\_training\_mini\_gpt.ipynb}
\reslink{Notebook}{03\_sampling\_strategies.ipynb}{../../notebooks/ch16/03\_sampling\_strategies.ipynb}
\reslink{Notebook}{04\_gemma\_generation.ipynb}{../../notebooks/ch16/04\_gemma\_generation.ipynb}
\reslink{Notebook}{05\_instruction\_tuning\_with\_lora.ipynb}{../../notebooks/ch16/05\_instruction\_tuning\_with\_lora.ipynb}
\reslink{Repo}{Author's official notebook}{https://github.com/fchollet/deep-learning-with-python-notebooks/blob/master/chapter16\_text-generation.ipynb}
\vskip 2pt
\end{frame}

\coursesection{01}{A brief history of sequence generation}{Twenty-five years from a recurrent cell to a consumer product.}

\framekicker{Section 16.1}
\begin{frame}[shrink=25]{This was a niche subtopic until very recently}
{\fontsize{9.4}{12.6}\selectfont\color{fg2}Generating sequences from a model only hit the mainstream around \textbf{2016}. The techniques are much older than that.\par}\vskip 4pt
\begin{center}
\includegraphics[width=0.92\linewidth,height=0.60\textheight,keepaspectratio]{figs/ch16-history.tex.pdf}
\end{center}
\vskip 3pt{\color{ink3}\fontsize{7.4}{9.4}\selectfont Sometimes a good idea takes fifteen years to get started.\par}
{\fontsize{9.4}{12.6}\selectfont\color{fg2}Douglas Eck applied LSTM to \textbf{music generation} in 2002, became a researcher at Google Brain, and in 2016 started Magenta. Alex Graves pioneered recurrent networks for new kinds of sequence data --- his 2013 work generating human-like \textbf{handwriting} from pen-position timeseries is seen by many as the turning point.\par}\vskip 4pt
\end{frame}

\framekicker{Section 16.1 $\cdot$ an aside worth keeping}
\begin{frame}[shrink=25]{The closest computers get to dreaming}
\begin{iquote}
\quotetext{\textit{"Generating sequential data is the closest computers get to dreaming."}}\quotecite{Alex Graves, in a commented-out line of a 2013 LaTeX file on arXiv}
\end{iquote}
{\fontsize{9.4}{12.6}\selectfont\color{fg2}The line was never meant to be read. Chollet cites it as a significant inspiration behind starting Keras --- which is to say that the framework this course is taught in has a remark hidden in a preprint somewhere in its lineage.\par}\vskip 4pt
{\fontsize{9.4}{12.6}\selectfont\color{fg2}Worth remembering when the field feels purely industrial. Much of what follows in this chapter came out of curiosity long before it came out of a product plan.\par}\vskip 4pt
\end{frame}

\framekicker{Section 16.1 $\cdot$ 2018}
\begin{frame}[shrink=25]{GPT combined three things, and invented none of them}
{\fontsize{9.4}{12.6}\selectfont\color{fg2}A year after \textit{Attention Is All You Need}, researchers at OpenAI published \textit{Improving Language Understanding by Generative Pre-Training}. They combined:\par}\vskip 4pt
\begin{center}
\includegraphics[width=0.92\linewidth,height=0.60\textheight,keepaspectratio]{figs/ch16-gpt-ingredients.tex.pdf}
\end{center}
\vskip 3pt{\color{ink3}\fontsize{7.4}{9.4}\selectfont Generative Pretrained Transformer --- three known ingredients, one surprising result.\par}
\begin{iband}[signal]
\textbf{GPT came with no modelling or training advancements.} What was interesting is that such a general setup beat more involved, task-specific techniques --- no complex text normalization, no per-benchmark architecture. Just \hilite{pretraining data and compute}.
\end{iband}
\end{frame}
\note{This is the slide to sit on for a professional audience. The lesson that generality plus scale beats specialised engineering is the single most consequential fact of the last decade of this field.}

\framekicker{Section 16.1 $\cdot$ four years}
\begin{frame}[shrink=25]{Then OpenAI scaled it with single-minded focus}
\begin{itable}{>{\raggedright\arraybackslash}p{0.188\textwidth}>{\raggedright\arraybackslash}p{0.1316\textwidth}>{\raggedright\arraybackslash}p{0.282\textwidth}>{\raggedright\arraybackslash}p{0.3384\textwidth}}
\thead{Model} & \thead{Year} & \thead{Parameters} & \thead{Training tokens} \\ \midrule
GPT-1 & 2018 & 117 million & 1 billion \\
GPT-2 & 2019 & 1.5 billion & more than 10 billion \\
GPT-3 & 2020 & 175 billion & around half a trillion \\
\end{itable}
{\fontsize{9.4}{12.6}\selectfont\color{fg2}\textbf{The architecture changed only slightly} across all three. What changed was the size of everything else --- and with each leap in scale, the quality of the generative output shot up substantially.\par}\vskip 4pt
{\fontsize{9.4}{12.6}\selectfont\color{fg2}That is roughly a \textbf{1,500-fold} increase in parameters and a \textbf{500-fold} increase in tokens over four years, for the same design.\par}\vskip 4pt
\end{frame}

\framekicker{Section 16.1 $\cdot$ what emerged at each scale}
\begin{frame}[shrink=25]{Three models, three different ways of being used}
\begin{minipage}[t]{0.3253\textwidth}
\begin{icard}
\cardh{GPT-1: fine-tune it}
Generative capability was a \textbf{by-product}. Evaluation was by attaching a dense classification head --- exactly what we did to RoBERTa in chapter 15.
\end{icard}
\end{minipage}%
\hfill
\begin{minipage}[t]{0.3253\textwidth}
\begin{icardaccent}
\cardh{GPT-2: show it examples}
Prompt with a handful of worked examples and get quality output \textbf{with no fine-tuning at all}. This is \textit{few-shot learning}.
\end{icardaccent}
\end{minipage}%
\hfill
\begin{minipage}[t]{0.3253\textwidth}
\begin{icardaccent}
\cardh{GPT-3: just describe it}
Examples were often unnecessary. A plain text description of the problem plus the input frequently sufficed.
\end{icardaccent}
\end{minipage}%
{\fontsize{9.4}{12.6}\selectfont\color{fg2}\textit{Few-shot learning} means teaching a model a new problem with only a handful of supervised examples --- \textbf{far too few for standard gradient descent}. Nothing is being trained; the capability was already in the pretrained distribution.\par}\vskip 4pt
\end{frame}

\framekicker{Section 16.1 $\cdot$ few-shot, concretely}
\begin{frame}[shrink=25]{What a GPT-2 prompt looked like}
{\fontsize{9.4}{12.6}\selectfont\color{fg2}To get a French translation of the word \textit{cheese}, you would prompt the model with a small pattern and let it continue:\par}\vskip 4pt
\ioutput{listings/ch16-01.txt}
{\fontsize{9.4}{12.6}\selectfont\color{fg2}By GPT-3, this was often enough:\par}\vskip 4pt
\ioutput{listings/ch16-02.txt}
{\fontsize{9.4}{12.6}\selectfont\color{fg2}Nothing in the model was updated between the two. The difference is entirely in \textbf{what the larger pretrained distribution already contains}.\par}\vskip 4pt
\end{frame}

\framekicker{Section 16.1 $\cdot$ the problems arrived with the capability}
\begin{frame}[shrink=25]{GPT-3's three unsolved problems, stated in 2020}
\begin{minipage}[t]{0.3253\textwidth}
\begin{icardbad}
\cardh{Hallucination}
Output veers from accurate to \textbf{completely false with zero indication}. Nothing in the output signals which is which.
\end{icardbad}
\end{minipage}%
\hfill
\begin{minipage}[t]{0.3253\textwidth}
\begin{icardbad}
\cardh{Prompt sensitivity}
Seemingly minor rewording triggers \textbf{large jumps} in performance, up or down. Behaviour is hard to predict.
\end{icardbad}
\end{minipage}%
\hfill
\begin{minipage}[t]{0.3253\textwidth}
\begin{icardbad}
\cardh{No adaptation}
The model cannot adapt to problems that were not \textbf{extensively featured} in its training data.
\end{icardbad}
\end{minipage}%
\begin{iband}[rose]
These were listed as open problems in 2020 and are \hilite{still open}. Everything later in this chapter --- instruction tuning, RLHF, RAG, chain-of-thought --- reduces their severity. \textbf{None of them removes the problem.}
\end{iband}
{\fontsize{9.4}{12.6}\selectfont\color{fg2}Nevertheless GPT-3's output was good enough to become the basis for ChatGPT, the first widespread consumer-facing generative model.\par}\vskip 4pt
\end{frame}
\note{For a professional audience this is the risk-register slide. Come back to it after RAG and after the reasoning section, and mark honestly which boxes have moved.}

\coursesection{02}{Training a mini-GPT}{Forty-one million parameters, a billion tokens, and the exact blueprint everyone else is using.}

\framekicker{Section 16.2}
\begin{frame}[shrink=25]{The data: 1\% of a colossal clean crawled corpus}
{\fontsize{9.4}{12.6}\selectfont\color{fg2}GPT-1 used \textbf{BooksCorpus} --- self-published books added without the explicit permission of their authors. It has since been taken down by its publishers.\par}\vskip 4pt
{\fontsize{9.4}{12.6}\selectfont\color{fg2}We use \textbf{C4}, the Colossal Clean Crawled Corpus released by Google in 2020. At 750 GB it is far more than we can train on, so we take under 1\%.\par}\vskip 4pt
\icode{python}{listing 16.1}{listings/ch16-03.py}
{\fontsize{9.4}{12.6}\selectfont\color{fg2}Fifty shards, each about 75 MB of raw text. One document per line, newlines escaped.\par}\vskip 4pt
\end{frame}

\framekicker{Section 16.2 $\cdot$ a warning worth reading twice}
\begin{frame}[shrink=25]{This is the most compute-intensive chapter in the book}
\begin{minipage}[t]{0.3253\textwidth}
\istat{\textasciitilde{}6 hours}{mini-GPT on a free Colab T4}
\end{minipage}%
\hfill
\begin{minipage}[t]{0.3253\textwidth}
\istat{\textasciitilde{}1 hour}{the whole chapter on an A100}
\end{minipage}%
\hfill
\begin{minipage}[t]{0.3253\textwidth}
\istat{1}{runtime restart needed mid-notebook}
\end{minipage}%
{\fontsize{9.4}{12.6}\selectfont\color{fg2}You will need to restart the Colab runtime partway through to free GPU memory before loading the larger pretrained model.\par}\vskip 4pt
\begin{iband}[amber]
You can always \textbf{read through} the expensive \inlinecode{fit()} calls and edit down the step count for quick experimentation. The learning is in the code, not in waiting for it.
\end{iband}
{\fontsize{9.4}{12.6}\selectfont\color{fg2}And if you are running this a few years from now, there is a good chance these examples are child's play to the hardware in front of you.\par}\vskip 4pt
\end{frame}

\framekicker{Section 16.2}
\begin{frame}[shrink=25]{What one document looks like}
{\fontsize{9.4}{12.6}\selectfont\color{fg2}As always, look at the data before designing anything around it.\par}\vskip 4pt
\ioutput{listings/ch16-04.txt}
{\fontsize{9.4}{12.6}\selectfont\color{fg2}Crawled web text: a headline, a newline, marketing copy. This is what the model's distribution will be made of, and it is worth holding that image when the output later looks like a \hilite{forum post that will not end}.\par}\vskip 4pt
\end{frame}

\framekicker{Section 16.2 $\cdot$ listing 16.2}
\begin{frame}[shrink=25]{SentencePiece: the same BPE, in C++}
{\fontsize{9.4}{12.6}\selectfont\color{fg2}Preprocessing this much text needs a fast tokenizer. The technique is exactly the byte-pair encoding we built in chapter 14; \textbf{SentencePiece} is a C++ implementation of it with a \inlinecode{detokenize()} that maps back to strings.\par}\vskip 4pt
\icode{python}{listing 16.2}{listings/ch16-05.py}
\ioutput{listings/ch16-06.txt}
{\fontsize{9.4}{12.6}\selectfont\color{fg2}A premade vocabulary of \textbf{32,000 terms}. KerasHub wraps the library as a Keras layer, so it composes with \inlinecode{tf.data}.\par}\vskip 4pt
\end{frame}

\framekicker{Section 16.2 $\cdot$ a design decision}
\begin{frame}[shrink=25]{GPT does not respect document boundaries}
{\fontsize{9.4}{12.6}\selectfont\color{fg2}When training GPT, the developers made \textbf{no attempt} to keep document boundaries from falling in the middle of a training sample. Documents are simply concatenated.\par}\vskip 4pt
{\fontsize{9.4}{12.6}\selectfont\color{fg2}Instead, a boundary is marked with a special token:\par}\vskip 4pt
\ioutput{listings/ch16-07.txt}
{\fontsize{9.4}{12.6}\selectfont\color{fg2}The model learns what a document boundary means from the token itself, rather than from how the data was cut. \textbf{Simplicity in the pipeline, paid for with one vocabulary entry.}\par}\vskip 4pt
\begin{iband}[signal]
This is a recurring pattern in LLM engineering: when a structural constraint is expensive to enforce in data, \hilite{encode it as a token and let the model learn it}.
\end{iband}
\end{frame}

\framekicker{Section 16.2 $\cdot$ listing 16.3}
\begin{frame}[shrink=25]{The pipeline, with interleaved shards}
{\fontsize{9.4}{12.6}\selectfont\color{fg2}Each shard is read independently and the outputs are interleaved, so every CPU core can tokenize a different file at the same time.\par}\vskip 4pt
\icode{python}{listing 16.3}{listings/ch16-08.py}
{\fontsize{9.4}{12.6}\selectfont\color{fg2}\inlinecode{rebatch(sequence\_length + 1)} windows the token stream into even 257-token samples; \inlinecode{x[:-1], x[1:]} is the offset-by-one split from chapter 15, unchanged.\par}\vskip 4pt
\end{frame}

\framekicker{Section 16.2 $\cdot$ the numbers}
\begin{frame}[shrink=25]{Just under a billion tokens}
\begin{minipage}[t]{0.3253\textwidth}
\istat{29,373}{batches}
\end{minipage}%
\hfill
\begin{minipage}[t]{0.3253\textwidth}
\istat{128 $\times$ 256}{samples $\times$ tokens per batch}
\end{minipage}%
\hfill
\begin{minipage}[t]{0.3253\textwidth}
\istat{\textasciitilde{}0.96 billion}{tokens of training data}
\end{minipage}%
\icode{python}{}{listings/ch16-09.py}
{\fontsize{9.4}{12.6}\selectfont\color{fg2}The batch count is not free to obtain: \inlinecode{ds.reduce(0, lambda c, \_: c + 1)} will count it, but tokenizing a dataset this size takes \textbf{several minutes on a fast CPU}. The \inlinecode{prefetch(8)} at the end of the pipeline is what keeps the GPU from waiting on it.\par}\vskip 4pt
\end{frame}

\framekicker{Section 16.2.1}
\begin{frame}[shrink=25]{GPT throws away the encoder}
{\fontsize{9.4}{12.6}\selectfont\color{fg2}The original GPT \textbf{simplifies} the sequence-to-sequence Transformer of chapter 15: no encoder, no cross-attention. Only the decoder stack remains, which means information can only travel left to right.\par}\vskip 4pt
\begin{center}
\includegraphics[width=0.92\linewidth,height=0.60\textheight,keepaspectratio]{figs/ch16-decoder-only.tex.pdf}
\end{center}
\vskip 3pt{\color{ink3}\fontsize{7.4}{9.4}\selectfont A question and its answer become one sequence, embedded by the same parameters.\par}
{\fontsize{9.4}{12.6}\selectfont\color{fg2}A decoder-only model can still handle sequence-to-sequence problems like question answering --- but the question and answer must be \textbf{combined into a single sequence}, and question tokens are treated no differently from answer tokens.\par}\vskip 4pt
\end{frame}

\framekicker{Section 16.2.1 $\cdot$ the trade}
\begin{frame}[shrink=25]{What the decoder-only bet costs, and what it buys}
\begin{minipage}[t]{0.494\textwidth}
\begin{icardwarn}
\cardh{The cost: less expressive inputs}
Given \textit{"Where is the capital of France?"}, the representation of \textbf{Where} cannot attend to \textbf{capital} or \textbf{France}. Even the input is processed one-directionally.
\end{icardwarn}
\end{minipage}%
\hfill
\begin{minipage}[t]{0.494\textwidth}
\begin{icardgood}
\cardh{The gain: no curation at all}
No need for datasets of input/output \textbf{pairs}. Everything is a single sequence, so you can train on any text on the internet, at any scale you can afford.
\end{icardgood}
\end{minipage}%
\begin{iband}[signal]
This is the trade that decided the field. Expressivity was sacrificed for \hilite{an unbounded supply of training data}, and the data won.
\end{iband}
\end{frame}
\note{Contrast explicitly with RoBERTa from chapter 15, which kept bidirectionality and is still the better choice when you only need to represent text rather than generate it.}

\framekicker{Section 16.2.1 $\cdot$ listing 16.4}
\begin{frame}[shrink=25]{The decoder block, minus cross-attention}
{\fontsize{9.4}{12.6}\selectfont\color{fg2}This is chapter 15's \inlinecode{TransformerDecoder} with the cross-attention layer removed and one addition: \textbf{dropout inside each block}.\par}\vskip 4pt
\icode{python}{listing 16.4}{listings/ch16-10.py}
{\fontsize{9.4}{12.6}\selectfont\color{fg2}Chapter 15 used a \textbf{single} Transformer layer, so one dropout at the end of the model sufficed. Here we stack \textbf{eight}, and regularisation inside each block matters.\par}\vskip 4pt
\end{frame}

\framekicker{Section 16.2.1 $\cdot$ listing 16.4}
\begin{frame}[shrink=25]{Its forward pass has only two stages now}
{\fontsize{9.4}{12.6}\selectfont\color{fg2}With cross-attention gone, \inlinecode{call()} takes one input and runs two stages instead of three.\par}\vskip 4pt
\icode{python}{listing 16.4 - call}{listings/ch16-11.py}
{\fontsize{9.4}{12.6}\selectfont\color{fg2}\inlinecode{use\_causal\_mask=True} is doing all the work that made the training objective valid in chapter 15. \textbf{It is the only thing standing between this model and reading its own labels.}\par}\vskip 4pt
\end{frame}

\framekicker{Section 16.2.1 $\cdot$ listing 16.5}
\begin{frame}[shrink=25]{Tying the input and output projections}
{\fontsize{9.4}{12.6}\selectfont\color{fg2}The two biggest weights in a Transformer both deal with vocabulary space: the token embedding, shape \textbf{(vocab\_size, hidden\_dim)}, and the output projection, shape \textbf{(hidden\_dim, vocab\_size)}.\par}\vskip 4pt
{\fontsize{9.4}{12.6}\selectfont\color{fg2}They are transposes of each other in shape. It turns out that making them transposes of each other in \textbf{value} is a good idea --- the final projection is a \textit{reverse embedding}, mapping hidden space back to token space.\par}\vskip 4pt
\begin{center}
\includegraphics[width=0.92\linewidth,height=0.60\textheight,keepaspectratio]{figs/ch16-tied.tex.pdf}
\end{center}
\vskip 3pt{\color{ink3}\fontsize{7.4}{9.4}\selectfont One matrix, used forward at the input and transposed at the output.\par}
\end{frame}

\framekicker{Section 16.2.1 $\cdot$ listing 16.5}
\begin{frame}[shrink=25]{A \inlinecode{reverse} argument is all it takes}
{\fontsize{9.4}{12.6}\selectfont\color{fg2}Chapter 15's \inlinecode{PositionalEmbedding}, with one branch added to \inlinecode{call()}.\par}\vskip 4pt
\icode{python}{listing 16.5}{listings/ch16-12.py}
{\fontsize{9.4}{12.6}\selectfont\color{fg2}With a 32,000-word vocabulary and 512 hidden dimensions, this saves \textbf{16 million parameters} --- a substantial share of a 41-million-parameter model.\par}\vskip 4pt
\end{frame}

\framekicker{Section 16.2.1 $\cdot$ listing 16.6}
\begin{frame}[shrink=25]{Eight blocks, and the model is done}
{\fontsize{9.4}{12.6}\selectfont\color{fg2}Everything assembled. Note the single \inlinecode{embedding} object used twice --- once forward, once reversed.\par}\vskip 4pt
\icode{python}{listing 16.6}{listings/ch16-13.py}
{\fontsize{9.4}{12.6}\selectfont\color{fg2}\textbf{41 million parameters} --- large for this book, and tiny against models today that range from a couple of billion to trillions. \inlinecode{mixed\_float16} trades some numerical fidelity for speed; chapter 18 explains it properly.\par}\vskip 4pt
\end{frame}

\framekicker{Section 16.2.2}
\begin{frame}[shrink=25]{Training a large Transformer is famously finicky}
{\fontsize{9.4}{12.6}\selectfont\color{fg2}The model is sensitive to parameter initialization and to the choice of optimizer. Stack many Transformer layers and it becomes easy to suffer \textbf{exploding gradients} --- parameters update too quickly and the loss never converges.\par}\vskip 4pt
{\fontsize{9.4}{12.6}\selectfont\color{fg2}A trick that works well: ease \textbf{linearly} into the full learning rate over a number of warmup steps, so the earliest updates are small.\par}\vskip 4pt
\icode{python}{listing 16.7}{listings/ch16-14.py}
{\fontsize{9.4}{12.6}\selectfont\color{fg2}A thousand steps of ramp, then a flat 2e-4. Plot it before training --- a schedule that is wrong is invisible in the loss curve until it is far too late.\par}\vskip 4pt
\end{frame}

\framekicker{Section 16.2.2 $\cdot$ listing 16.8}
\begin{frame}[shrink=25]{One pass over a billion tokens, in eight epochs}
{\fontsize{9.4}{12.6}\selectfont\color{fg2}Eight epochs, not because the data repeats, but so we can check validation loss and accuracy periodically during a single pass.\par}\vskip 4pt
\icode{python}{listing 16.8}{listings/ch16-15.py}
{\fontsize{9.4}{12.6}\selectfont\color{fg2}We use \textbf{3$\times$ fewer parameters} than GPT-1 and \textbf{100$\times$ fewer training steps}. Even so, this is the most computationally expensive training run in the entire book.\par}\vskip 4pt
\end{frame}

\framekicker{Section 16.2.2 $\cdot$ sidebar}
\begin{frame}[shrink=25]{What is a logit?}
{\fontsize{9.4}{12.6}\selectfont\color{fg2}The output projection has \textbf{no softmax activation}. Its outputs are unnormalized log probabilities --- exponentiate each and normalise them to sum to 1, which is all softmax does, and you have probabilities.\par}\vskip 4pt
{\fontsize{9.4}{12.6}\selectfont\color{fg2}The common term for an unnormalized log probability is a \textbf{logit}, and logits are easier to work with when generating text --- as the sampling section will show.\par}\vskip 4pt
\begin{minipage}[t]{0.494\textwidth}
\begin{icard}
\cardh{Softmax in the model}
\inlinecode{Dense(n, activation="softmax")} plus \inlinecode{SparseCategoricalCrossentropy()}. The model outputs probabilities.
\end{icard}
\end{minipage}%
\hfill
\begin{minipage}[t]{0.494\textwidth}
\begin{icardaccent}
\cardh{Softmax in the loss}
\inlinecode{Dense(n)} plus \inlinecode{SparseCategoricalCrossentropy(from\_logits=True)}. The model outputs logits --- \textbf{more numerically stable, and easier to sample from}.
\end{icardaccent}
\end{minipage}%
\end{frame}

\framekicker{Section 16.2.2}
\begin{frame}[shrink=25]{36\% --- and the model is deliberately undertrained}
{\fontsize{9.4}{12.6}\selectfont\color{fg2}After training, the model predicts the next token correctly about \textbf{36\%} of the time on the validation set. As always, next-token accuracy is a crude heuristic rather than a measure of usefulness.\par}\vskip 4pt
\begin{iband}[amber]
The validation loss is \textbf{still ticking down} after every epoch. This model is not converged and we know it --- unsurprising with a hundred times fewer training steps than GPT-1. Training longer would help; it would also cost time and money.
\end{iband}
{\fontsize{9.4}{12.6}\selectfont\color{fg2}Being explicit about an undertrained model is itself a discipline. Most published curves stop where the budget ran out, not where the model stopped improving.\par}\vskip 4pt
\end{frame}

\coursesection{03}{Decoding and sampling strategies}{The model gives you a distribution. What you do with it is a separate design.}

\framekicker{Section 16.2.3 $\cdot$ listing 16.9}
\begin{frame}[shrink=25]{The naive generation loop}
{\fontsize{9.4}{12.6}\selectfont\color{fg2}The same approach as the Shakespeare generator: feed a prompt, take the prediction at the last position, append, repeat.\par}\vskip 4pt
\icode{python}{listing 16.9}{listings/ch16-16.py}
{\fontsize{9.4}{12.6}\selectfont\color{fg2}Correct, readable --- and it takes \textbf{minutes} to produce 64 tokens. That is the puzzle to solve before anything else.\par}\vskip 4pt
\end{frame}

\framekicker{Section 16.2.3}
\begin{frame}[shrink=25]{What it produces, and why it is slow}
\ioutput{listings/ch16-17.txt}
{\fontsize{9.4}{12.6}\selectfont\color{fg2}Two separate problems, and it is worth naming them apart. \textbf{One:} it repeats itself. \textbf{Two:} it took minutes, when training ran at 200,000 tokens per second on the same hardware.\par}\vskip 4pt
{\fontsize{9.4}{12.6}\selectfont\color{fg2}The second is a compilation problem. \inlinecode{fit()} and \inlinecode{predict()} compile the per-batch computation --- \inlinecode{keras.ops} calls are lifted out of Python and optimised by the backend. Calling the model \textbf{directly} runs the forward pass live and \hilite{unoptimized at every step}.\par}\vskip 4pt
\end{frame}

\framekicker{Section 16.2.3 $\cdot$ listing 16.10}
\begin{frame}[shrink=25]{Pad the input so the shape never changes}
{\fontsize{9.4}{12.6}\selectfont\color{fg2}\inlinecode{predict()} compiles for a \textbf{specific input shape}. Our sequence grows by one token each step, which would trigger recompilation every call. Pad to full length up front and the shape is constant.\par}\vskip 4pt
\icode{python}{listing 16.10}{listings/ch16-18.py}
\ioutput{listings/ch16-19.txt}
{\fontsize{9.4}{12.6}\selectfont\color{fg2}From minutes to \textbf{under half a second}. Same maths, same weights, same output --- a pure engineering win.\par}\vskip 4pt
\end{frame}
\note{Worth flagging for anyone deploying: a large share of real-world inference cost is lost to recompilation from variable shapes. It rarely shows up as an error, only as a bill.}

\framekicker{Section 16.2.3 $\cdot$ sidebar}
\begin{frame}[shrink=25]{One inefficiency remains. Can you spot it?}
{\fontsize{9.4}{12.6}\selectfont\color{fg2}Each call runs the model over the \textbf{entire} sequence, then throws away everything but one position. The sequence changes by a single token between steps.\par}\vskip 4pt
{\fontsize{9.4}{12.6}\selectfont\color{fg2}With the RNN in chapter 15 we kept the state and computed one token at a time. A causal Transformer has an equivalent notion of state --- you just have to look for it.\par}\vskip 4pt
\begin{center}
\includegraphics[width=0.92\linewidth,height=0.60\textheight,keepaspectratio]{figs/ch16-kvcache.tex.pdf}
\end{center}
\vskip 3pt{\color{ink3}\fontsize{7.4}{9.4}\selectfont Attention is the only place information crosses positions, and past keys and values never change.\par}
\end{frame}

\framekicker{Section 16.2.3 $\cdot$ sidebar}
\begin{frame}[shrink=25]{Key-value caching, and why every real system has it}
{\fontsize{9.4}{12.6}\selectfont\color{fg2}For a query at a given position, attention scores come from dotting that query with \textbf{all previous key vectors}, and the output combines \textbf{all previous value vectors}. Past input is fixed, and the causal mask forbids looking ahead --- so those vectors never change.\par}\vskip 4pt
{\fontsize{9.4}{12.6}\selectfont\color{fg2}Cache every key and value at every layer, and you have the Transformer equivalent of an RNN's state.\par}\vskip 4pt
\begin{iband}[amber]
Model input goes from \textit{as long as the output} to \textbf{one token}. On a sequence thousands of tokens long this is a \hilite{thousandfold speed-up}. Any efficient implementation of generative sampling includes it.
\end{iband}
{\fontsize{9.4}{12.6}\selectfont\color{fg2}Implementing it is clunky --- saving and reusing intermediate arrays from every attention layer --- which is exactly why you should use a library that has already done it.\par}\vskip 4pt
\end{frame}

\framekicker{Section 16.2.4}
\begin{frame}[shrink=25]{The repetition is not a bug}
{\fontsize{9.4}{12.6}\selectfont\color{fg2}The model repeats \textit{"get a sense of what you are doing"} over and over. This follows \textbf{directly from the training objective}.\par}\vskip 4pt
{\fontsize{9.4}{12.6}\selectfont\color{fg2}The model predicts the most likely next token across a billion words on many topics. Where there is no obvious continuation, an effective strategy is to guess \textbf{common words or repeated patterns} --- and the model learns this almost immediately.\par}\vskip 4pt
\begin{iband}[signal]
Stop training very early and the model would generate the word \inlinecode{"the"} incessantly, because \textit{the} is the most common word in English. Repetition is the objective \hilite{working correctly}, being asked the wrong question.
\end{iband}
\end{frame}

\framekicker{Section 16.2.4}
\begin{frame}[shrink=25]{The output is a distribution, not a token}
{\fontsize{9.4}{12.6}\selectfont\color{fg2}Throughout the generative loop we have taken the most likely token. But the output is not a token --- it is a probability distribution across all \textbf{32,000} entries in the vocabulary.\par}\vskip 4pt
{\fontsize{9.4}{12.6}\selectfont\color{fg2}Taking the most likely output at each step is called \textbf{greedy search}. It is the most straightforward use of the predictions, and hardly the only one. Adding randomness lets us explore the learned distribution more broadly, and keeps us out of high-probability loops.\par}\vskip 4pt
\begin{center}
\includegraphics[width=0.92\linewidth,height=0.60\textheight,keepaspectratio]{figs/ch16-sampling.tex.pdf}
\end{center}
\vskip 3pt{\color{ink3}\fontsize{7.4}{9.4}\selectfont Figure 16.3 --- three strategies over the same distribution.\par}
\end{frame}

\framekicker{Section 16.2.4}
\begin{frame}[shrink=25]{Refactor: make the strategy a parameter}
{\fontsize{9.4}{12.6}\selectfont\color{fg2}Pass a function that maps model predictions to a chosen token. Everything that follows is a different \inlinecode{sample\_fn}.\par}\vskip 4pt
\icode{python}{}{listings/ch16-20.py}
{\fontsize{9.4}{12.6}\selectfont\color{fg2}The sampling strategy is now a \textbf{first-class control} rather than a hardcoded \inlinecode{argmax} buried in a loop --- which is how every serving framework exposes it too.\par}\vskip 4pt
\end{frame}

\framekicker{Section 16.2.4 $\cdot$ random sampling}
\begin{frame}[shrink=25]{Sample the distribution directly}
{\fontsize{9.4}{12.6}\selectfont\color{fg2}\inlinecode{keras.random.categorical()} passes predictions through a softmax and samples the resulting distribution.\par}\vskip 4pt
\icode{python}{}{listings/ch16-21.py}
\ioutput{listings/ch16-22.txt}
{\fontsize{9.4}{12.6}\selectfont\color{fg2}More diverse, and no longer stuck in loops --- but now it \textbf{explores too much}. The output jumps around without continuity. We have traded one failure for its opposite.\par}\vskip 4pt
\end{frame}

\framekicker{Section 16.2.4 $\cdot$ temperature}
\begin{frame}[shrink=25]{One number that sharpens or flattens the distribution}
{\fontsize{9.4}{12.6}\selectfont\color{fg2}Dividing the logits before the softmax rescales how peaked the distribution is.\par}\vskip 4pt
\begin{center}
\includegraphics[width=0.92\linewidth,height=0.60\textheight,keepaspectratio]{figs/ch16-temperature.tex.pdf}
\end{center}
\vskip 3pt{\color{ink3}\fontsize{7.4}{9.4}\selectfont Temperature acts on the logits, before the softmax --- not on the probabilities after it.\par}
\begin{ibullets}
\item \textbf{Low temperature} makes all logits larger before the softmax, so the most likely output becomes even more likely.
\item \textbf{High temperature} makes them smaller, so the distribution spreads out.
\end{ibullets}
\end{frame}

\framekicker{Section 16.2.4 $\cdot$ temperature, observed}
\begin{frame}[shrink=25]{T = 2.0: not English any more}
{\fontsize{9.4}{12.6}\selectfont\color{fg2}The failure at high temperature is unmistakable.\par}\vskip 4pt
\ioutput{listings/ch16-23.txt}
{\fontsize{9.4}{12.6}\selectfont\color{fg2}Subword fragments, stray identifiers, tokens from other languages. At high temperature the distribution is flat enough that rare tokens win regularly, and \textbf{rare tokens in a 32,000 vocabulary are mostly debris}.\par}\vskip 4pt
\end{frame}

\framekicker{Section 16.2.4 $\cdot$ temperature, observed}
\begin{frame}[shrink=25]{T = 0.8 and T = 0.2: fluency against repetition}
\ioutput{listings/ch16-24.txt}
\ioutput{listings/ch16-25.txt}
{\fontsize{9.4}{12.6}\selectfont\color{fg2}At low temperature the behaviour converges on \textbf{greedy search}, repeating patterns. Temperature is a dial between two known failure modes, not a fix for either.\par}\vskip 4pt
\end{frame}

\framekicker{Section 16.2.4 $\cdot$ top-K}
\begin{frame}[shrink=25]{Restrict the choice to the K most likely tokens}
{\fontsize{9.4}{12.6}\selectfont\color{fg2}A different shaping technique: keep only the top \textit{K} candidates and sample among them. \inlinecode{keras.ops.top\_k} does the selection.\par}\vskip 4pt
\icode{python}{}{listings/ch16-26.py}
\begin{iband}[signal]
\textbf{Temperature and top-K are not the same knob.} A low temperature makes likely tokens more likely but rules \hilite{nothing} out. Top-K sets the probability of everything outside the K candidates to \hilite{zero}
\end{iband}
\end{frame}

\framekicker{Section 16.2.4 $\cdot$ top-K, observed}
\begin{frame}[shrink=25]{K = 5, K = 20, and the two combined}
\ioutput{listings/ch16-27.txt}
\ioutput{listings/ch16-28.txt}
{\fontsize{9.4}{12.6}\selectfont\color{fg2}The two compose: sampling the top five candidates at a temperature of 0.5 is a common production default.\par}\vskip 4pt
\end{frame}

\framekicker{Section 16.2.4 $\cdot$ what to take away}
\begin{frame}[shrink=25]{Sampling is a control, and there are more of them}
\begin{itable}{>{\raggedright\arraybackslash}p{0.2068\textwidth}>{\raggedright\arraybackslash}p{0.4136\textwidth}>{\raggedright\arraybackslash}p{0.3196\textwidth}}
\thead{Strategy} & \thead{What it does} & \thead{Failure mode} \\ \midrule
\textbf{Greedy} & Takes the argmax at every step & Repeats phrases indefinitely \\
\textbf{Random} & Samples the full categorical distribution & Wanders, no continuity \\
\textbf{Temperature} & Scales logits before the softmax & A dial between the two above \\
\textbf{Top-K} & Zeroes everything outside K candidates & K too small collapses to greedy \\
\textbf{Beam search} & Keeps several candidate chains alive at once & Expensive; can still be bland \\
\end{itable}
{\fontsize{9.4}{12.6}\selectfont\color{fg2}With top-K our mini-GPT produces plausible English --- with \textbf{little apparent utility}. That matches the GPT-1 result exactly: the generated output was a curiosity, and the state-of-the-art numbers came from fine-tuned classifiers.\par}\vskip 4pt
\end{frame}

\framekicker{Section 16.2.4 $\cdot$ the gap to today}
\begin{frame}[shrink=25]{What separates this from a modern LLM}
\begin{minipage}[t]{0.494\textwidth}
\istat{100$\times$}{more parameters needed}
\end{minipage}%
\hfill
\begin{minipage}[t]{0.494\textwidth}
\istat{1,000$\times$}{more training steps needed}
\end{minipage}%
\leadpar{And that is the entire gap. \textbf{The training recipe we just used is the exact blueprint everyone training LLMs uses today.}}
{\fontsize{9.4}{12.6}\selectfont\color{fg2}The only missing pieces are a very large compute budget and some tricks for training across multiple machines --- which chapter 18 covers. If we spent it, we would see the same leaps in quality OpenAI observed.\par}\vskip 4pt
\end{frame}
\note{This is the slide that reframes the whole chapter for a professional audience: the barrier to building a frontier model is capital and infrastructure, not undisclosed algorithms.}

\coursesection{04}{Using a pretrained LLM}{A billion parameters, an instruction dataset, and a memory problem.}

\framekicker{Section 16.3}
\begin{frame}[shrink=25]{Why almost nobody pretrains}
{\fontsize{9.4}{12.6}\selectfont\color{fg2}Given how prohibitively expensive pretraining is, most of the industry has centred on using models from a \textbf{short list of companies}. The concern is not only cost.\par}\vskip 4pt
\begin{minipage}[t]{0.3253\textwidth}
\istat{1.3 million kWh}{estimated to train Llama 2}
\end{minipage}%
\hfill
\begin{minipage}[t]{0.3253\textwidth}
\istat{45,000}{American households' daily power}
\end{minipage}%
\hfill
\begin{minipage}[t]{0.3253\textwidth}
\istat{smaller than GPT-3}{and this is the cheap case}
\end{minipage}%
\begin{iband}[rose]
Generative model training is now a \textbf{large percentage of total data-centre power consumption} at big technology companies. If every organisation using an LLM pretrained its own, the energy use would be a noticeable percentage of \hilite{global consumption}.
\end{iband}
\end{frame}

\framekicker{Section 16.3.1 $\cdot$ listing 16.11}
\begin{frame}[shrink=25]{Loading Gemma 3, one billion parameters}
{\fontsize{9.4}{12.6}\selectfont\color{fg2}We use the smallest Gemma 3 variant --- almost exactly a billion parameters, trained on roughly \textbf{2 trillion tokens}, two thousand times more than our mini-GPT.\par}\vskip 4pt
\icode{python}{listing 16.11}{listings/ch16-29.py}
{\fontsize{9.4}{12.6}\selectfont\color{fg2}\inlinecode{CausalLM} is a high-level task API, like the \inlinecode{ImageClassifier} and \inlinecode{ImageSegmenter} of earlier chapters. It combines a tokenizer and a correctly initialised architecture into one Keras model, and loads matching weights.\par}\vskip 4pt
\begin{iband}[amber]
You must accept the \textbf{Gemma Terms of Use} on Kaggle before the weights will download, and authenticate with \inlinecode{kagglehub.login()}. The terms prohibit uses such as generating spam or hate speech --- licence terms like this are becoming standard.
\end{iband}
\end{frame}

\framekicker{Section 16.3.1}
\begin{frame}[shrink=25]{Where a billion parameters actually sit}
\ioutput{listings/ch16-30.txt}
{\fontsize{9.4}{12.6}\selectfont\color{fg2}Two things worth reading closely. The vocabulary is \textbf{262,144} terms, eight times ours. And \inlinecode{ReversibleEmbedding} is the \textbf{tied embedding} we built by hand --- 302 million of the billion parameters, counted once because input and output share them.\par}\vskip 4pt
\end{frame}

\framekicker{Section 16.3.1}
\begin{frame}[shrink=25]{It is a fancy autocomplete for the internet}
{\fontsize{9.4}{12.6}\selectfont\color{fg2}\inlinecode{CausalLM} provides \inlinecode{generate()}, compilable with the sampling strategies we just built.\par}\vskip 4pt
\ioutput{listings/ch16-31.txt}
{\fontsize{9.4}{12.6}\selectfont\color{fg2}Far more coherent than mini-GPT --- and \textbf{still not useful}. Asked how to make brownies, it wrote the \textit{forum post asking the question}, because that is a likely continuation in crawled web text.\par}\vskip 4pt
\end{frame}

\framekicker{Section 16.3.1 $\cdot$ prompting}
\begin{frame}[shrink=25]{Change the prompt, change which part of the distribution you visit}
{\fontsize{9.4}{12.6}\selectfont\color{fg2}One way to fix the output is to prompt with a longer input that makes the desired continuation the likely one.\par}\vskip 4pt
\ioutput{listings/ch16-32.txt}
\begin{iband}[signal]
It is tempting to imagine the model \textit{interpreting} the prompt conversationally. \textbf{Nothing of the sort is going on.} We constructed a prompt for which an actual recipe is a more likely continuation than someone asking for help.
\end{iband}
\end{frame}

\framekicker{Section 16.3.1 $\cdot$ the honest assessment}
\begin{frame}[shrink=25]{Prompt engineering is useful and hard to control}
{\fontsize{9.4}{12.6}\selectfont\color{fg2}You can go much further: prompt with natural-language instructions about the role the model should fill (\textit{"You are a large language model that gives short, helpful answers"}), or with a long list of topics to avoid.\par}\vskip 4pt
\begin{iquote}
\quotetext{If this all sounds a bit hand-wavy and hard to control, \textbf{that's a good assessment.} Attempting to visit different parts of a model's distribution through prompting is often useful, but predicting how a model will respond to a given prompt is very difficult.}\quotecite{Section 16.3.1}
\end{iquote}
{\fontsize{9.4}{12.6}\selectfont\color{fg2}Treat prompts as \textbf{empirical artefacts}: version them, test them against a fixed evaluation set, and expect them to break when the model behind them changes.\par}\vskip 4pt
\end{frame}

\framekicker{Section 16.3.1 $\cdot$ hallucination}
\begin{frame}[shrink=25]{There is always a most-likely next token}
{\fontsize{9.4}{12.6}\selectfont\color{fg2}A model will always say \textbf{something}. Finding regions of the distribution with no grounding in fact is easy:\par}\vskip 4pt
\ioutput{listings/ch16-33.txt}
{\fontsize{9.4}{12.6}\selectfont\color{fg2}Utter nonsense --- and the model could not find a \textbf{more likely} way to complete the prompt. There is no mechanism by which it could decline; declining is itself just another continuation, and one this model was never trained to prefer.\par}\vskip 4pt
\begin{iband}[rose]
Hallucination and uncontrollable output are \textbf{fundamental} problems with language models. If there is a silver bullet, \hilite{we have yet to find it}.
\end{iband}
\end{frame}

\framekicker{Section 16.3.2}
\begin{frame}[shrink=25]{Instruction fine-tuning: bend the output, do not relearn the language}
{\fontsize{9.4}{12.6}\selectfont\color{fg2}Feed the model input/output pairs --- a user instruction followed by a response --- combined into a single sequence, with markers so it is clear where each ends.\par}\vskip 4pt
{\fontsize{9.4}{12.6}\selectfont\color{fg2}Then train with the \textbf{same next-token loss} used for pretraining. The precise markup does not matter, as long as it is consistent.\par}\vskip 4pt
\begin{iband}[signal]
We are not learning a latent space for language here --- that was done over trillions of tokens. We are \hilite{nudging the learned representation} to control the tone and content of the output.
\end{iband}
\end{frame}

\framekicker{Section 16.3.2 $\cdot$ listing 16.12}
\begin{frame}[shrink=25]{Dolly: 15,000 handwritten instruction-response pairs}
{\fontsize{9.4}{12.6}\selectfont\color{fg2}Databricks employees contributed 15,000 instructions with handwritten responses.\par}\vskip 4pt
\icode{python}{listing 16.12}{listings/ch16-34.py}
{\fontsize{9.4}{12.6}\selectfont\color{fg2}Examples carrying extra \textbf{context} are discarded here --- those are the RAG-shaped ones, which return later.\par}\vskip 4pt
\end{frame}

\framekicker{Section 16.3.2}
\begin{frame}[shrink=25]{The template gives the examples a predictable shape}
\ioutput{listings/ch16-35.txt}
{\fontsize{9.4}{12.6}\selectfont\color{fg2}Gemma is not a sequence-to-sequence model like our translator. But by training on prompts with this structure and generating only what follows the \inlinecode{[response]} marker, we can \textbf{use it in a sequence-to-sequence setting}.\par}\vskip 4pt
\icode{python}{}{listings/ch16-36.py}
{\fontsize{9.4}{12.6}\selectfont\color{fg2}Batch size \textbf{2}. That number is about to become important.\par}\vskip 4pt
\end{frame}

\framekicker{Section 16.3.2}
\begin{frame}[shrink=25]{Inside the preprocessor}
{\fontsize{9.4}{12.6}\selectfont\color{fg2}\inlinecode{CausalLM} wraps two objects: a \textbf{preprocessor} layer and a \textbf{backbone} model. Preprocessing is automatic inside \inlinecode{fit()} and \inlinecode{predict()}, but running it by hand shows what it does.\par}\vskip 4pt
\ioutput{listings/ch16-37.txt}
{\fontsize{9.4}{12.6}\selectfont\color{fg2}The padding mask tracks which inputs are padded zeros. The \inlinecode{sample\_weight} tensor restricts the loss to \textbf{response tokens} --- we do not care about loss on a fixed user prompt, and certainly not on padding.\par}\vskip 4pt
\end{frame}

\framekicker{Section 16.3.2}
\begin{frame}[shrink=25]{It is the same language-model setup underneath}
{\fontsize{9.4}{12.6}\selectfont\color{fg2}Print the token IDs beside the labels and the offset-by-one pattern from chapter 15 is visible unchanged.\par}\vskip 4pt
\ioutput{listings/ch16-38.txt}
{\fontsize{9.4}{12.6}\selectfont\color{fg2}Each label is the next token value. \textbf{Instruction tuning is not a new objective} --- it is the pretraining objective, run on carefully chosen data with the loss masked to the parts we care about.\par}\vskip 4pt
\end{frame}

\framekicker{Section 16.3.3}
\begin{frame}[shrink=25]{Why fit() would run out of memory right now}
{\fontsize{9.4}{12.6}\selectfont\color{fg2}We loaded the model and ran generation without trouble. So why would training crash on a 16 GB GPU?\par}\vskip 4pt
\begin{itable}{>{\raggedright\arraybackslash}p{0.5264\textwidth}>{\raggedright\arraybackslash}p{0.4136\textwidth}}
\thead{What has to be in memory} & \thead{Size} \\ \midrule
Model weights & \textbf{3.7 GB} \\
Adam: gradients, velocity, momentum --- 3 floats per parameter & \textbf{\textasciitilde{}11 GB} \\
Intermediate values in the forward pass & a few GB \\
\textbf{Total} & \textbf{over 16 GB} \\
\end{itable}
\begin{iband}[amber]
This is the common shape of LLM training: because parameter counts are so large, \textbf{GPU throughput is a secondary concern to fitting the model in accelerator memory at all}.
\end{iband}
\end{frame}

\framekicker{Section 16.3.3}
\begin{frame}[shrink=25]{Freezing saves optimizer state, not just gradients}
{\fontsize{9.4}{12.6}\selectfont\color{fg2}Earlier chapters froze parts of a model during fine-tuning. What went unmentioned: \textbf{frozen parameters need no optimizer variables at all}, because they never update. That is where the memory goes.\par}\vskip 4pt
{\fontsize{9.4}{12.6}\selectfont\color{fg2}Researchers have experimented extensively with which parameters to leave unfrozen in a Transformer. The answer, perhaps intuitively, is the ones in the \textbf{attention mechanism}.\par}\vskip 4pt
{\fontsize{9.4}{12.6}\selectfont\color{fg2}But the attention layers still hold hundreds of millions of parameters. \hilite{Can we do better?}\par}\vskip 4pt
\end{frame}

\framekicker{Section 16.3.3 $\cdot$ LoRA}
\begin{frame}[shrink=25]{Freeze the kernel, learn a low-rank correction}
{\fontsize{9.4}{12.6}\selectfont\color{fg2}In 2021 researchers at Microsoft proposed \textbf{LoRA} --- Low-Rank Adaptation --- specifically for this memory problem. Start from an ordinary linear projection:\par}\vskip 4pt
\icode{python}{}{listings/ch16-39.py}
{\fontsize{9.4}{12.6}\selectfont\color{fg2}LoRA freezes \inlinecode{kernel} and adds a \textbf{low-rank decomposition} of the update alongside it: two matrices projecting down to an inner rank and back out.\par}\vskip 4pt
\end{frame}

\framekicker{Section 16.3.3 $\cdot$ LoRA}
\begin{frame}[shrink=25]{Alpha and beta, and a rank in between}
{\fontsize{9.4}{12.6}\selectfont\color{fg2}Three weights instead of one.\par}\vskip 4pt
\icode{python}{}{listings/ch16-40.py}
\begin{center}
\includegraphics[width=0.92\linewidth,height=0.60\textheight,keepaspectratio]{figs/ch16-lora.tex.pdf}
\end{center}
\vskip 3pt{\color{ink3}\fontsize{7.4}{9.4}\selectfont Figure 16.4 --- the low-rank decomposition holds far fewer parameters than the kernel it corrects.\par}
\end{frame}

\framekicker{Section 16.3.3 $\cdot$ the arithmetic}
\begin{frame}[shrink=25]{Four million parameters become thirty-two thousand}
\begin{minipage}[t]{0.3253\textwidth}
\istat{2048 $\times$ 2048}{kernel shape}
\end{minipage}%
\hfill
\begin{minipage}[t]{0.3253\textwidth}
\istat{4,194,304}{frozen parameters}
\end{minipage}%
\hfill
\begin{minipage}[t]{0.3253\textwidth}
\istat{32,768}{trainable, at rank 8}
\end{minipage}%
{\fontsize{9.4}{12.6}\selectfont\color{fg2}The update does \textbf{not} have the expressive power of the original kernel --- at the narrow middle, the entire update must pass through eight floats.\par}\vskip 4pt
\begin{iband}[signal]
That is the bet, and it holds: \textbf{during fine-tuning you no longer need the expressive power you needed during pretraining.} The representation is already built; you are only steering it.
\end{iband}
\end{frame}

\framekicker{Section 16.3.3 $\cdot$ listing 16.13}
\begin{frame}[shrink=25]{In KerasHub it is one line}
{\fontsize{9.4}{12.6}\selectfont\color{fg2}The LoRA authors suggest freezing the entire Transformer and adding LoRA weights only to the \textbf{query and key} projections in attention.\par}\vskip 4pt
\icode{python}{listing 16.13}{listings/ch16-41.py}
{\fontsize{9.4}{12.6}\selectfont\color{fg2}The same thing written out layer by layer, which is what you would edit to add trainable parameters elsewhere:\par}\vskip 4pt
\icode{python}{the verbose equivalent}{listings/ch16-42.py}
{\fontsize{9.4}{12.6}\selectfont\color{fg2}Written out this way you can see the two decisions the one-liner makes for you: \textbf{which layers} get LoRA, and at \textbf{what rank}. Adding it to the value projection, or to only the later blocks, is a one-line change from here.\par}\vskip 4pt
\end{frame}

\framekicker{Section 16.3.3}
\begin{frame}[shrink=25]{A thousandfold cut in trainable parameters}
\ioutput{listings/ch16-43.txt}
{\fontsize{9.4}{12.6}\selectfont\color{fg2}The weights still occupy 3.7 GB. But the \textbf{trainable} parameters are now 5 MB --- which takes the optimizer state from many gigabytes down to megabytes.\par}\vskip 4pt
{\fontsize{9.4}{12.6}\selectfont\color{fg2}Note that total parameters went \textit{up} slightly, from 999.9 M to 1,001.2 M: the alpha and beta matrices are new weights. \textbf{LoRA adds parameters to the model and removes them from the optimizer} --- and the optimizer was the problem.\par}\vskip 4pt
\end{frame}

\framekicker{Section 16.3.3 $\cdot$ listing 16.14}
\begin{frame}[shrink=25]{Now the fine-tuning fits}
{\fontsize{9.4}{12.6}\selectfont\color{fg2}One epoch over roughly 7,000 instruction pairs, at the same small learning rate discipline used for RoBERTa in chapter 15.\par}\vskip 4pt
\icode{python}{listing 16.14}{listings/ch16-44.py}
{\fontsize{9.4}{12.6}\selectfont\color{fg2}We reach \textbf{55\%} accuracy at guessing the next word of the response, against 36\% for mini-GPT. \inlinecode{weighted\_metrics} is what makes the \inlinecode{sample\_weight} mask reach the metric, so the score is over response tokens only.\par}\vskip 4pt
\end{frame}

\framekicker{Section 16.3.3 $\cdot$ the result}
\begin{frame}[shrink=25]{It answers questions now}
\ioutput{listings/ch16-45.txt}
{\fontsize{9.4}{12.6}\selectfont\color{fg2}Much better. The model now \textbf{responds} to a question rather than carrying on the thought of the prompt text --- and it emits \inlinecode{[end]} when it is finished, because the training data taught it where responses stop.\par}\vskip 4pt
\end{frame}

\framekicker{Section 16.3.3 $\cdot$ the honest test}
\begin{frame}[shrink=25]{Have we solved hallucination? Not at all.}
\ioutput{listings/ch16-46.txt}
{\fontsize{9.4}{12.6}\selectfont\color{fg2}Identical nonsense, now delivered in a helpful tone. \textbf{Instruction tuning changed the format of the answer, not its relationship to truth.}\par}\vskip 4pt
\begin{iband}[amber]
One thing that helps: train on many instruction/response pairs where the desired response is \textit{"I don't know"} or \textit{"As a language model, I cannot help you with that"}. This teaches the model to avoid whole topics where it would answer badly --- \hilite{a behaviour, not an understanding}.
\end{iband}
\end{frame}
\note{For a professional audience this is the slide that should govern any deployment decision. Fluency and helpfulness are trainable. Correctness is not, by this mechanism.}

\coursesection{05}{Going further with LLMs}{RLHF, multimodality, retrieval, and models that show their working.}

\framekicker{Section 16.4.1}
\begin{frame}[shrink=25]{Handwritten data has a ceiling, and it is human}
{\fontsize{9.4}{12.6}\selectfont\color{fg2}What we just did is called \textbf{supervised fine-tuning}: we curated, by hand, a list of prompts and responses we wanted.\par}\vskip 4pt
\begin{ibullets}
\item Manually written examples are \textbf{slow and expensive} to come by, and will almost always become the bottleneck.
\item The approach is limited by the \textbf{human performance ceiling} on the task. To beat human performance, we cannot supervise with human-written output.
\end{ibullets}
{\fontsize{9.4}{12.6}\selectfont\color{fg2}The real thing to optimise is our \textbf{preference} for some responses over others. With a large enough sample of people this is perfectly well defined --- the difficulty is translating \textit{our preferences} into a loss function. That is what RLHF attempts.\par}\vskip 4pt
\end{frame}

\framekicker{Section 16.4.1}
\begin{frame}[shrink=25]{RLHF in five steps}
\begin{center}
\includegraphics[width=0.92\linewidth,height=0.60\textheight,keepaspectratio]{figs/ch16-rlhf.tex.pdf}
\end{center}
\vskip 3pt{\color{ink3}\fontsize{7.4}{9.4}\selectfont The reward model is a proxy for human preference --- a smaller Transformer that reads a sequence and outputs one float.\par}
{\fontsize{9.4}{12.6}\selectfont\color{fg2}Responses to be ranked may be handwritten, generated by the model, or even written by other chatbots. Collecting rankings is expensive and slow --- but \textbf{still faster than writing every desired response by hand}.\par}\vskip 4pt
\end{frame}

\framekicker{Section 16.4.1 $\cdot$ listing 16.15}
\begin{frame}[shrink=25]{The simplest possible RLHF algorithm}
{\fontsize{9.4}{12.6}\selectfont\color{fg2}Do not be intimidated by \textit{reinforcement learning}: it refers to any setup where a model learns by making predictions (\textbf{actions}) and receiving feedback (\textbf{rewards}). The model's own predictions become its training data.\par}\vskip 4pt
\icode{python}{listing 16.15}{listings/ch16-47.py}
{\fontsize{9.4}{12.6}\selectfont\color{fg2}In practice you would \textbf{not discard} the bad responses --- a bad response is a good signal about what not to do --- and you would use specialised gradient updates over all responses. OpenAI used this setup to go from GPT-3's pretrained parameters to the first ChatGPT.\par}\vskip 4pt
\end{frame}

\framekicker{Section 16.4.1 $\cdot$ listing 16.16}
\begin{frame}[shrink=25]{An instruction-tuned checkpoint, four times larger}
{\fontsize{9.4}{12.6}\selectfont\color{fg2}Companies releasing pretrained models normally also release specialised \textbf{instruction-tuned} versions. Let us load one --- 4 billion parameters, the largest model in this book.\par}\vskip 4pt
\icode{python}{listing 16.16}{listings/ch16-48.py}
\begin{minipage}[t]{0.494\textwidth}
\begin{icard}
\cardh{float32 for training}
4 bytes per parameter. 32-bit floats represent very small values, which helps keep training \textbf{gradients stable}.
\end{icard}
\end{minipage}%
\hfill
\begin{minipage}[t]{0.494\textwidth}
\begin{icardaccent}
\cardh{bfloat16 for inference}
2 bytes per parameter. No training here, so no gradient-stability concern --- and \textbf{many gigabytes saved}.
\end{icardaccent}
\end{minipage}%
\end{frame}

\framekicker{Section 16.4.1}
\begin{frame}[shrink=25]{The template must match what the tuning used}
{\fontsize{9.4}{12.6}\selectfont\color{fg2}This checkpoint has its own prompt format. Again the exact text does not matter --- \textbf{matching} it does.\par}\vskip 4pt
\icode{python}{}{listings/ch16-49.py}
\ioutput{listings/ch16-50.txt}
{\fontsize{9.4}{12.6}\selectfont\color{fg2}Pretrained on \textbf{14 trillion tokens}, then extensively tuned --- supervised fine-tuning, RLHF, and distillation from a larger \textit{teacher} model.\par}\vskip 4pt
\end{frame}

\framekicker{Section 16.4.1 $\cdot$ the same trick question}
\begin{frame}[shrink=25]{It refuses the bait --- and then invents the date}
\ioutput{listings/ch16-51.txt}
{\fontsize{9.4}{12.6}\selectfont\color{fg2}This is \textbf{not a new modelling technique}. It is the result of extensive training on trick questions like this one, with responses like this one.\par}\vskip 4pt
\begin{iband}[rose]
And look closely: having refused to hallucinate a president, the model \textbf{made up today's date}. Removing hallucinations is \hilite{whack-a-mole}, and this example is the demonstration.
\end{iband}
\end{frame}

\framekicker{Section 16.4.2}
\begin{frame}[shrink=25]{Multimodality is a sequence problem}
\leadpar{The Transformer is not a text-specific model. It is a highly effective model for learning patterns in \textbf{sequence data}.}
{\fontsize{9.4}{12.6}\selectfont\color{fg2}So if we can coerce another data type into a sequence representation, we can feed it to a Transformer and train on it. That is the whole idea.\par}\vskip 4pt
\begin{center}
\includegraphics[width=0.92\linewidth,height=0.60\textheight,keepaspectratio]{figs/ch16-multimodal.tex.pdf}
\end{center}
\vskip 3pt{\color{ink3}\fontsize{7.4}{9.4}\selectfont Figure 16.6 --- hard text tokens and soft image tokens spliced into one sequence.\par}
\end{frame}

\framekicker{Section 16.4.2 $\cdot$ hard and soft tokens}
\begin{frame}[shrink=25]{An image becomes 256 tokens that were never in a vocabulary}
{\fontsize{9.4}{12.6}\selectfont\color{fg2}The Gemma model we just loaded carries a separate \textbf{420-million-parameter image encoder}. It cuts an input image into 256 patches and encodes each as a vector of the same dimensionality as Gemma's hidden state.\par}\vskip 4pt
\begin{itable}{>{\raggedright\arraybackslash}p{0.2444\textwidth}>{\raggedright\arraybackslash}p{0.3478\textwidth}>{\raggedright\arraybackslash}p{0.3478\textwidth}}
\thead{} & \thead{Hard tokens} & \thead{Soft tokens} \\ \midrule
Source & A token ID & The vision encoder's output \\
Possible values & One row of the embedding matrix & \textbf{Any vector at all} \\
Count here & As many as the text has & 256 per image \\
Loss & Trained to predict them & Loss is \textbf{zeroed} at these positions \\
\end{itable}
{\fontsize{9.4}{12.6}\selectfont\color{fg2}Each image embeds as a \textbf{(256, 2560)} sequence, spliced into the text sequence after the token embedding layer.\par}\vskip 4pt
\end{frame}

\framekicker{Section 16.4.2}
\begin{frame}[shrink=25]{Asking questions about a picture}
{\fontsize{9.4}{12.6}\selectfont\color{fg2}The prompt is now a dictionary with two keys, and the preprocessor is told how many images to expect.\par}\vskip 4pt
\icode{python}{}{listings/ch16-52.py}
\ioutput{listings/ch16-53.txt}
{\fontsize{9.4}{12.6}\selectfont\color{fg2}The special token \inlinecode{<start\_of\_image>} is expanded into 256 placeholder positions, which are then \textbf{replaced} by the soft tokens from the vision encoder.\par}\vskip 4pt
\end{frame}

\framekicker{Section 16.4.2 $\cdot$ sidebar}
\begin{frame}[shrink=25]{Foundation models}
{\fontsize{9.4}{12.6}\selectfont\color{fg2}Once a model handles images and audio, \textit{large language model} becomes misleading. The umbrella term is \textbf{foundation model}: any model trained on broad data, generally with self-supervision at scale, that can be fine-tuned to a wide range of downstream tasks.\par}\vskip 4pt
\begin{center}
\includegraphics[width=0.92\linewidth,height=0.60\textheight,keepaspectratio]{figs/ch16-foundation.tex.pdf}
\end{center}
\vskip 3pt{\color{ink3}\fontsize{7.4}{9.4}\selectfont The hallmarks: a self-supervised reconstruction loss, and no specialisation to a single task.\par}
{\fontsize{9.4}{12.6}\selectfont\color{fg2}This is a \textbf{striking and recent shift}. Rather than training from scratch on your own dataset, you are often better off getting a rich representation from a foundation model and specialising it --- with the downside of running billions of parameters, which is \hilite{hardly a fit for every application}.\par}\vskip 4pt
\end{frame}

\framekicker{Section 16.4.3}
\begin{frame}[shrink=25]{Two reasons an LLM is not a search engine}
\begin{minipage}[t]{0.494\textwidth}
\begin{icardbad}
\cardh{It makes things up}
False \textit{facts} that were not in the training data but could be interpolated from it. This ranges from \textbf{misleading to dangerous}.
\end{icardbad}
\end{minipage}%
\hfill
\begin{minipage}[t]{0.494\textwidth}
\begin{icardbad}
\cardh{Its knowledge has a cutoff}
At best, the date it was pretrained. Training is expensive and cannot run continuously, so knowledge of the world just \textbf{stops}.
\end{icardbad}
\end{minipage}%
{\fontsize{9.4}{12.6}\selectfont\color{fg2}Nobody wants a search engine that can only tell you about things that happened six months ago. But if we treat the LLM as \hilite{conversational software} that can handle any sequence data in a prompt, we can use it as the \textbf{interface} to information retrieved by more traditional search.\par}\vskip 4pt
\end{frame}

\framekicker{Section 16.4.3 $\cdot$ RAG}
\begin{frame}[shrink=25]{Retrieve first, then generate}
{\fontsize{9.4}{12.6}\selectfont\color{fg2}RAG takes a user question, runs a query to pull in additional text context --- from a database, a search engine, anything --- and puts that context \textbf{straight into the prompt}.\par}\vskip 4pt
\ioutput{listings/ch16-54.txt}
\begin{center}
\includegraphics[width=0.92\linewidth,height=0.60\textheight,keepaspectratio]{figs/ch16-rag.tex.pdf}
\end{center}
\vskip 3pt{\color{ink3}\fontsize{7.4}{9.4}\selectfont A vector database keys documents by their embeddings and returns the nearest ones to the query embedding.\par}
\end{frame}

\framekicker{Section 16.4.3 $\cdot$ what RAG buys}
\begin{frame}[shrink=25]{Three problems it addresses, and one it does not}
\begin{isteps}
\item It gives an obvious way to \textbf{work around the cutoff date} of the model.
\item It lets the model access \textbf{private data} --- a company can use a publicly-trained LLM as an interface to information it holds internally.
\item It helps \textbf{factually ground} the model: an LLM is much less likely to invent facts on a topic when correct context sits in the prompt.
\end{isteps}
\begin{iband}[rose]
There is \textbf{no silver bullet that stops hallucination entirely.} RAG makes it less likely on topics you retrieved context for. It does nothing for topics you did not.
\end{iband}
{\fontsize{9.4}{12.6}\selectfont\color{fg2}One pleasing detail: the vector database's \textit{query}, \textit{key}, and \textit{value} vocabulary is not an accident --- attention borrowed those terms \textbf{from database systems} in the first place.\par}\vskip 4pt
\end{frame}

\framekicker{Section 16.4.4}
\begin{frame}[shrink=25]{Grade-school maths defied progress for years}
{\fontsize{9.4}{12.6}\selectfont\color{fg2}For years LLMs were abysmal at maths problems and logic puzzles. A model might answer a problem from its training data perfectly --- substitute a few names or numbers and it became evident it had \textbf{no grasp of what it was solving}.\par}\vskip 4pt
{\fontsize{9.4}{12.6}\selectfont\color{fg2}For most NLP problems the recipe was easy: more data, better benchmark score. Grade-school word problems did not respond to it.\par}\vskip 4pt
\begin{iquote}
\quotetext{In 2023 researchers at Google noticed that prompting the model with a few examples of \textit{showing your work} made the model do the same --- and that by attending to its own written-out steps, it reached correct solutions far more often. They called it \textbf{chain-of-thought prompting}.}\quotecite{Section 16.4.4}
\end{iquote}
{\fontsize{9.4}{12.6}\selectfont\color{fg2}Another group found the examples were not even necessary: prompting with \textit{"Let's think step by step"} was enough.\par}\vskip 4pt
\end{frame}

\framekicker{Section 16.4.4 $\cdot$ training it in}
\begin{frame}[shrink=25]{Chain-of-thought fine-tuning, in five steps}
{\fontsize{9.4}{12.6}\selectfont\color{fg2}Models like OpenAI's o1 and DeepSeek's r1 made headlines by training a model to \textit{think out loud}. The approach is very close to RLHF.\par}\vskip 4pt
\begin{isteps}
\item Collect basic maths and reasoning problems with their desired answers.
\item Generate, \textbf{with randomness}, a number of responses to each.
\item Find responses with a correct answer by \textbf{string parsing} --- prompt the model to mark its final answer with a specific token.
\item Run supervised fine-tuning on the correct responses, \textbf{including all the intermediate output}.
\item Repeat.
\end{isteps}
{\fontsize{9.4}{12.6}\selectfont\color{fg2}The answer check acts as the \textbf{environment}; the generated outputs are the \textbf{actions}. In practice you would use a more complex gradient step that learns from incorrect responses too.\par}\vskip 4pt
\end{frame}

\framekicker{Section 16.4.4 $\cdot$ observed}
\begin{frame}[shrink=25]{The same prompt, twice, with random sampling}
{\fontsize{9.4}{12.6}\selectfont\color{fg2}\textit{Judy wrote a 2-page letter to 3 friends twice a week for 3 months. How many letters did she write?}\par}\vskip 4pt
\ioutput{listings/ch16-55.txt}
\ioutput{listings/ch16-56.txt}
{\fontsize{9.4}{12.6}\selectfont\color{fg2}The first attempt got hung up on the \textbf{superfluous detail} that each letter has two pages. The second is right. Same model, same prompt, different sample --- which is exactly why step 3 of the recipe is \textit{check the answer}.\par}\vskip 4pt
\end{frame}

\framekicker{Section 16.4.4}
\begin{frame}[shrink=25]{Where else verifiable answers exist}
\begin{center}
\includegraphics[width=0.92\linewidth,height=0.60\textheight,keepaspectratio]{figs/ch16-cot.tex.pdf}
\end{center}
\vskip 3pt{\color{ink3}\fontsize{7.4}{9.4}\selectfont The loop only closes where correctness can be checked automatically.\par}
{\fontsize{9.4}{12.6}\selectfont\color{fg2}The same idea applies wherever a text prompt has an obvious, \textbf{verifiable} answer. \textbf{Coding} is the important one: prompt the model for code, then actually run it to test the quality of the response.\par}\vskip 4pt
{\fontsize{9.4}{12.6}\selectfont\color{fg2}One trend is clear across all these domains: as a model learns harder questions, it spends \textbf{more and more time showing its work} before answering. Think of it as learning to search over its own potential solutions.\par}\vskip 4pt
\end{frame}

\coursesection{06}{Where are LLMs heading next?}{More parameters is the obvious answer. It is not quite the right one.}

\framekicker{Section 16.5}
\begin{frame}[shrink=25]{A fixed budget buys flops, and you must choose how to spend them}
{\fontsize{9.4}{12.6}\selectfont\color{fg2}Given a fixed pretraining budget --- say a million dollars --- you are buying a fixed number of floating-point operations. You can spend them on \textbf{more data}, or on a \textbf{bigger model}.\par}\vskip 4pt
\begin{center}
\includegraphics[width=0.92\linewidth,height=0.60\textheight,keepaspectratio]{figs/ch16-scaling.tex.pdf}
\end{center}
\vskip 3pt{\color{ink3}\fontsize{7.4}{9.4}\selectfont Recent research found GPT-3 was far too big for the compute it was given.\par}
{\fontsize{9.4}{12.6}\selectfont\color{fg2}Training a \textbf{smaller model on more data} would have produced better performance. So model sizes have trended flatter while data sizes have trended up.\par}\vskip 4pt
\end{frame}

\framekicker{Section 16.5}
\begin{frame}[shrink=25]{Two reasons parameter counts stopped climbing}
\begin{minipage}[t]{0.494\textwidth}
\begin{icardaccent}
\cardh{GPT-3 was undertrained}
175 billion parameters against roughly half a trillion tokens. We now know the balance was wrong.
\end{icardaccent}
\end{minipage}%
\hfill
\begin{minipage}[t]{0.494\textwidth}
\begin{icardaccent}
\cardh{Deployment cost}
It is often worth sacrificing performance for a smaller model that fits cheaper hardware. \textbf{A really good model does not help if it is prohibitively expensive to run.}
\end{icardaccent}
\end{minipage}%
{\fontsize{9.4}{12.6}\selectfont\color{fg2}This does \textbf{not} mean scaling has stopped. More compute does generally give better performance, and we have yet to see any sign of an asymptote where next-token prediction levels off. Billions of dollars continue to go into finding out what emerges.\par}\vskip 4pt
\end{frame}

\framekicker{Section 16.5 $\cdot$ the supply problem}
\begin{frame}[shrink=25]{We are starting to run out of pretraining data}
{\fontsize{9.4}{12.6}\selectfont\color{fg2}Technology companies are having trouble finding more high-quality, public, \textbf{human-written} content to throw at pretraining.\par}\vskip 4pt
\begin{iband}[rose]
Models are starting to \textit{eat their own tail} --- training on a significant portion of content created by other LLMs, which brings a whole other set of concerns.
\end{iband}
{\fontsize{9.4}{12.6}\selectfont\color{fg2}This is one reason reinforcement learning is drawing so much attention. \textbf{If you can build a difficult, self-contained environment that generates new problems, you can keep training on the model's own output} --- with no need to scrounge the web for more quality text.\par}\vskip 4pt
\end{frame}

\framekicker{Section 16.5 $\cdot$ the fundamental problem}
\begin{frame}[shrink=25]{None of this is a silver bullet}
\begin{iquote}
\quotetext{At the end of the day, the fundamental problem remains that LLMs are \textbf{wildly inefficient at learning compared to humans}. Model capabilities only come from training on many orders of magnitude more text than people will read in their lifetimes.}\quotecite{Section 16.5}
\end{iquote}
{\fontsize{9.4}{12.6}\selectfont\color{fg2}As scaling continues, so will more fundamental research into models that learn quickly from limited data.\par}\vskip 4pt
{\fontsize{9.4}{12.6}\selectfont\color{fg2}Still --- LLMs represent the ability to build \textbf{fluent natural language interfaces}, and that alone brings a massive shift in what can be accomplished with computing devices. This chapter laid out the basic recipe by which they get there.\par}\vskip 4pt
\end{frame}

\framekicker{Common failure modes}
\begin{frame}[shrink=25]{Four ways LLM work goes wrong in practice}
\begin{minipage}[t]{0.494\textwidth}
\begin{icardbad}
\cardh{Uncompiled or reshaping generation}
Calling the model directly, or letting the input shape change each step, costs \textbf{two to three orders of magnitude} in speed. No error is raised --- only a bill.
\end{icardbad}
\end{minipage}%
\hfill
\begin{minipage}[t]{0.494\textwidth}
\begin{icardbad}
\cardh{Prompt template mismatch}
An instruction-tuned checkpoint expects the exact template it was tuned with. Use a different one and quality drops sharply, silently.
\end{icardbad}
\end{minipage}%


\vskip 4pt

\begin{minipage}[t]{0.494\textwidth}
\begin{icardwarn}
\cardh{Full fine-tuning on a small GPU}
Adam needs three extra floats per parameter. The model loads and generates fine, then \inlinecode{fit()} dies on the first step. Use LoRA.
\end{icardwarn}
\end{minipage}%
\hfill
\begin{minipage}[t]{0.494\textwidth}
\begin{icardwarn}
\cardh{Mistaking fluency for correctness}
Instruction tuning and RLHF change the \textbf{form} of an answer. Neither makes it true. Ground with retrieval, and verify where you can.
\end{icardwarn}
\end{minipage}%
\end{frame}

\framekicker{Summary}
\begin{frame}[shrink=25]{What has to survive this chapter (1 of 2)}
\begin{isteps}
\item \textbf{An LLM is three ingredients}: the Transformer architecture, a language modelling task, and a large amount of unlabelled text.
\item \textbf{GPT is decoder-only} --- expressivity traded for an unbounded supply of training data, and the data won.
\item \textbf{Tied embeddings} use one matrix forward at the input and transposed at the output, saving a large share of the parameters.
\item \textbf{Warmup} matters: large Transformers are sensitive to early updates, and exploding gradients are the common failure.
\item \textbf{Sampling is a separate design} from the model. Greedy repeats, random wanders, temperature dials between them, top-K rules tokens out.
\item \textbf{Key-value caching} turns generation from quadratic re-computation into one token per step.
\end{isteps}
\end{frame}

\framekicker{Summary}
\begin{frame}[shrink=25]{What has to survive this chapter (2 of 2)}
\begin{isteps}
\item \textbf{Instruction fine-tuning} is the pretraining objective on curated data with the loss masked to responses --- it bends output, it does not relearn language.
\item \textbf{LoRA} freezes the kernel and learns a low-rank correction, cutting trainable parameters a thousandfold and optimizer memory with them.
\item \textbf{RLHF} replaces handwritten answers with ranked preferences and a reward model, which lifts the human performance ceiling.
\item \textbf{Multimodality is a sequence problem} --- encode images as soft tokens and splice them into the text sequence.
\item \textbf{RAG} puts retrieved context in the prompt, addressing the knowledge cutoff and private data, and reducing but not removing hallucination.
\item \textbf{All LLMs hallucinate.} Fluency, helpfulness, and refusal are trainable behaviours; truth is not one of them.
\end{isteps}
\vskip 4pt
\reslink{NOTEBOOK}{05\_instruction\_tuning\_with\_lora.ipynb}{../../course-slides/notebooks/ch16/05\_instruction\_tuning\_with\_lora.ipynb}
\reslink{PAPER}{Hu et al., LoRA (2021)}{https://arxiv.org/abs/2106.09685}
\reslink{NEXT}{Chapter 17 --- Image generation}{../ch17/index.html}
\vskip 2pt
\end{frame}

\end{document}
